\section{Introduction}
\label{sec:introduction}

Automated recognition of crackles and wheezes can reduce the manual effort
involved in respiratory sound assessment\cite{hsu2021hflung}.
Training these classifiers requires expert annotations that are costly
to obtain at scale\cite{zhang2024opera}. Existing respiratory
datasets\cite{rocha2019icbhi,zhang2022sprsound} offer an opportunity to
expand supervision by reusing labeled recordings. Their annotations
describe related acoustic attributes but differ in granularity and
prediction targets, so their task labels are not directly interchangeable.

To reuse annotations that are not directly interchangeable, a common
strategy is to map source labels to a shared prediction target, as in
SPRSound's fusion experiments and LungMix's cross-dataset
transfer\cite{zhang2022sprsound,ge2025lungmix}. This enables reuse under an
aligned task definition, but coarsening labels can discard distinctions
needed by another task. Representing sound composition offers a more
flexible way to use label information: Multi-breath and PC-MCL employ
multilabel supervision within ICBHI\cite{chua2024multibreath,jeong2026pcmcl}.
Across datasets, DCASE combines class mappings and masked supervision while
retaining dataset-specific evaluation within a shared temporal sound-event
detection system\cite{cornell2024dcase}. These studies provide ways to align
labels and retain their information, but our setting requires sharing
supervision across distinct cycle-level and event-level prediction targets.
For example, SPRSound's Wheeze annotations contain information relevant to
ICBHI's four-class cycle task, although its binary event task requires only
normal/abnormal predictions\cite{zhang2022sprsound}. The available supervision
is therefore richer than the required output. We aim to jointly learn from
this information while preserving each dataset's native prediction target.

To this end, we propose LSAA (Learning Shared Acoustic Attributes).
A shared encoder learns acoustic representations through abnormality,
crackle, and wheeze supervision from ICBHI cycles and SPRSound events.
Each annotation supervises only the attributes it supports.
Task-specific readouts map the shared attribute predictions to four-class
cycle and binary event outputs, allowing fine supervision to guide
learning while preserving different native tasks.

In this work, we demonstrate:
\begin{enumerate}
\setlength{\itemsep}{1pt}
\setlength{\parsep}{0pt}
\item We propose LSAA, a joint learning framework that connects heterogeneous
respiratory sound annotations through shared acoustic attributes in one model.
\item We develop task-specific readouts over shared abnormality, crackle,
and wheeze predictions, enabling fine-grained supervision to support
distinct native outputs.
\item We evaluate native-task performance on ICBHI and SPRSound,
investigate attribute transfer and auxiliary supervision on HF
Lung\cite{hsu2021hflung}, and assess external generalization on
KAUH\cite{fraiwan2021kauh}.
\end{enumerate}
